\documentclass{IEEEtran}
\usepackage{fancyvrb}
\fvset{numbers=left,fontsize=\fontsize{8}{10}\selectfont,numbersep=6pt}

\usepackage{amsmath}
\usepackage{float}
\usepackage{hyperref}

\newcommand{\pannelwidth}{0.32}
\setlength{\parskip}{1ex}
\setlength{\baselineskip}{0.1ex}

\title{BMI3 ICA mini-project proposal title}
\author{Group X}

\date{\today}

\begin{document}

\maketitle
\begin{abstract}
    Summary of proposal (~300 words): Your proposal in a few sentences. Including Objectives, Significance, and a brief description of the proposed Work. A comment on anticipated outcomes (deliverables).
\end{abstract}

\section{Survey of Background Literature}

\subsection{Overview}

\IEEEPARstart{G}{ive} Context and a brief argument for the need for progress in this research area and what you propose to do, in simple, non-technical terms (for the benefit of non-technical experts). Also, briefly emphasize the impact and outcomes from this work in tangible terms, such as societal benefits.

\subsection{Background}

Give a thorough introduction to the background science relevant to the general area of the proposed research (much of this comes from the review of the topic). Explain any relevant theory and standard techniques used in the area, along with any acronyms and jargon typical to the field. Also, provide background material on any new methodologies you intend to introduce to the field. Remember that it is always better to avoid the use of abbreviations and acronyms as much as possible.

\subsection{Relevance/Impact}

In the context provided by the background science, layout the argument that justifies the need for research in this area. Point out the significance of these works to the field and practice. Comment on its potential impact on the development of the scientific area and the society/sector as a whole. What are we going to learn? Why is it important to learn this? What would be the biological significance of solving the problem?

\section{Proposed Methodology}

Outline the proposed algorithm(s) specific to the proposed research. Justify the choice of algorithms as opposed to alternatives. Avoid giving too much detailed information, just outline the general approach and why you chose to use it.

\section{Research Plan}

State the long- and short-term objectives of your research program. Outline specific projects planned to meet these goals, including the timelines for completion of each stage, methods to be used, and dissemination of results. 

\section{Resources}

Provide details on the dataset you plan to use.

\section{References}

Add your used references. Any (consistent) format is fine.

\vspace{2cm}

[NOTE: be brief, strictly no more than 2 pages.]

\end{document}